% ======================================================================
% Sigles et acronymes : source unique pour le rapport, l'entente et la présentation.
% Syntaxe (sans le % initial) : \newacronym[options acro] {clé} {forme courte} {forme longue}
%   (espaces ajoutés ici pour que LaTeX Workshop ignore cette ligne d'exemple)
%   options : foreign = terme anglais d'origine (en italique, entre parenthèses)
%             sort    = clé de tri, utile si la forme courte est accentuée
% \newacronym est un alias de \DeclareAcronym (paquet acro), défini dans
% shared/macros.tex : ce n'est PAS la commande du paquet glossaries.
% Pourquoi : LaTeX Workshop ne lit que la syntaxe \newacronym pour proposer
% les clés en autocomplétion après \ac{, \acs{, \acl{.
% La liste (\listesigles) est triée automatiquement : l'ordre ici est libre.
% ======================================================================
\newacronym{CER}{CER}{comité d'éthique de la recherche}
\newacronym{CGIPN}{CGIPN}{Centre de gouvernance de l'information des Premières Nations}
\newacronym[foreign = differential privacy]{DP}{DP}{confidentialité différentielle}
\newacronym[sort = EDI]{EDI}{ÉDI}{équité, diversité et inclusion}
\newacronym{EPTC}{EPTC~2}{Énoncé de politique des trois Conseils : Éthique de la recherche avec des êtres humains}
\newacronym[foreign = Federated Averaging]{FedAvg}{FedAvg}{moyenne fédérée}
\newacronym[foreign = federated learning]{FL}{FL}{apprentissage fédéré}
\newacronym{IA}{IA}{intelligence artificielle}
% PCAP® : marque déposée du CGIPN, le symbole ® est demandé (vérifié).
\newacronym{PCAP}{PCAP\marquedeposee}{propriété, contrôle, accès et possession}
\newacronym{PI}{PI}{propriété intellectuelle}
\newacronym{UQO}{UQO}{Université du Québec en Outaouais}
